% LaTeX Template for AI Problem Analysis Output (v2.0)
% AMC12B 2023 Problem 10 Strategic Analysis
% Generated using AMC12_COMC_Problem_Solving_Prompt.md framework

\documentclass[12pt]{article}
\usepackage[utf8]{inputenc}
\usepackage{amsmath, amssymb, amsthm}
\usepackage{geometry}
\usepackage{xcolor}
\usepackage[most]{tcolorbox}
\usepackage{enumitem}
\usepackage{fancyhdr}

% Page setup
\geometry{margin=1in}
\setlength{\headheight}{14.5pt}
\pagestyle{fancy}
\fancyhf{}
\rhead{AMC12B 2023 Problem 10 Analysis}
\lhead{\today}

% Color definitions
\definecolor{problemconfig}{RGB}{230, 242, 255}    % Light blue
\definecolor{strategic}{RGB}{255, 230, 242}       % Light pink
\definecolor{tactical}{RGB}{242, 255, 230}        % Light green
\definecolor{techniques}{RGB}{255, 242, 230}      % Light yellow
\definecolor{verification}{RGB}{255, 230, 230}    % Light red
\definecolor{solution}{RGB}{230, 255, 255}        % Light cyan
\definecolor{competition}{RGB}{242, 242, 255}     % Light purple
\definecolor{proofwriting}{RGB}{240, 230, 255}    % Light lavender
\definecolor{gold}{RGB}{218, 165, 32}

% Box styles
\newtcolorbox{problemconfigbox}{
    colback=problemconfig,
    colframe=blue,
    title=Problem Configuration Analysis,
    fonttitle=\bfseries,
    arc=3pt,
    boxrule=1pt,
    breakable
}

\newtcolorbox{strategicbox}{
    colback=strategic,
    colframe=purple,
    title=Strategic Framework Application,
    fonttitle=\bfseries,
    arc=3pt,
    boxrule=1pt,
    breakable
}

\newtcolorbox{tacticalbox}{
    colback=tactical,
    colframe=green,
    title=Tactical Methodology Selection,
    fonttitle=\bfseries,
    arc=3pt,
    boxrule=1pt,
    breakable
}

\newtcolorbox{techniquesbox}{
    colback=techniques,
    colframe=orange,
    title=Conversion Techniques Application,
    fonttitle=\bfseries,
    arc=3pt,
    boxrule=1pt,
    breakable
}

\newtcolorbox{verificationbox}{
    colback=verification,
    colframe=red,
    title=Verification Strategy Design,
    fonttitle=\bfseries,
    arc=3pt,
    boxrule=1pt,
    breakable
}

\newtcolorbox{solutionbox}{
    colback=solution,
    colframe=cyan,
    title=Solution Roadmap,
    fonttitle=\bfseries,
    arc=3pt,
    boxrule=1pt,
    breakable
}

\newtcolorbox{competitionbox}{
    colback=competition,
    colframe=purple,
    title=Competition Strategy Integration,
    fonttitle=\bfseries,
    arc=3pt,
    boxrule=1pt,
    breakable
}

\newtcolorbox{keyinsightbox}{
    colback=yellow!20,
    colframe=gold,
    title=Key Insight,
    fonttitle=\bfseries,
    arc=3pt,
    boxrule=2pt,
    leftrule=5pt,
    breakable
}

\newtcolorbox{warningbox}{
    colback=red!10,
    colframe=red!50,
    title=Warning: Common Pitfall,
    fonttitle=\bfseries,
    arc=3pt,
    boxrule=2pt,
    leftrule=5pt,
    breakable
}

\newtcolorbox{tipbox}{
    colback=green!10,
    colframe=green!50,
    title=Pro Tip,
    fonttitle=\bfseries,
    arc=3pt,
    boxrule=2pt,
    leftrule=5pt,
    breakable
}

\begin{document}

\title{AMC12B 2023 Problem 10\\Strategic Problem Analysis}
\author{Competition Math Framework v2.1}
\date{\today}
\maketitle

\section*{Problem Statement}

In the $xy$-plane, a circle of radius $4$ with center on the positive $x$-axis is tangent to the $y$-axis at the origin, and a circle with radius $10$ with center on the positive $y$-axis is tangent to the $x$-axis at the origin. What is the slope of the line passing through the two points at which these circles intersect?

$$\textbf{(A)}\ \dfrac{2}{7} \qquad\textbf{(B)}\ \dfrac{3}{7} \qquad\textbf{(C)}\ \dfrac{2}{\sqrt{29}} \qquad\textbf{(D)}\ \dfrac{1}{\sqrt{29}} \qquad\textbf{(E)}\ \dfrac{2}{5}$$

\begin{problemconfigbox}
\subsection*{A. Problem Classification}
\begin{itemize}[leftmargin=*]
    \item \textbf{Problem Type}: To Find (specific slope value)
    \item \textbf{Competition Context}: AMC12B 2023, Problem 10
    \item \textbf{Difficulty Band}: Mid-band (P10) - Intermediate level, 3-4 minutes expected
    \item \textbf{Mathematical Domain}: Geometry (Circle Geometry, Coordinate Geometry)
    \item \textbf{Estimated Time}: 3-4 minutes with verification
\end{itemize}

\subsection*{B. Problem Structure Identification}
\begin{itemize}[leftmargin=*]
    \item \textbf{Given Information}: 
    \begin{itemize}
        \item Circle 1: radius $r_1 = 4$, center on positive $x$-axis, tangent to $y$-axis at origin
        \item Circle 2: radius $r_2 = 10$, center on positive $y$-axis, tangent to $x$-axis at origin
        \item Two circles intersect at two points
    \end{itemize}
    \item \textbf{Constraints}: 
    \begin{itemize}
        \item Centers must be on positive coordinate axes
        \item Tangency condition determines exact center locations
        \item Circles must intersect (not just touch externally or contain one another)
    \end{itemize}
    \item \textbf{Objective}: Find the slope of the line through the two intersection points
    \item \textbf{Hidden Assumptions}: 
    \begin{itemize}
        \item The tangency at origin means the circles touch the axes at $(0,0)$
        \item For a circle tangent to an axis, the distance from center to axis equals the radius
        \item The two intersection points are distinct
    \end{itemize}
    \item \textbf{Answer Choice Leverage}: 
    \begin{itemize}
        \item Answer choices include simple fractions and expressions with radicals
        \item Options (A) and (B) have denominator 7, suggesting possible relationship
        \item Options (C) and (D) involve $\sqrt{29}$, suggesting distance or Pythagorean result
        \item Option (E) is the simplest form with denominator 5
        \item Can eliminate choices that don't match geometric constraints
    \end{itemize}
\end{itemize}

\subsection*{C. Strategic Orientation}
\begin{itemize}[leftmargin=*]
    \item \textbf{Hypothesis}: Given two specific circles with tangency conditions
    \item \textbf{Conclusion}: Determine slope of chord connecting intersection points
    \item \textbf{Notation}: 
    \begin{itemize}
        \item Center of first circle: $C_1 = (4, 0)$ (since tangent to $y$-axis at origin with radius 4)
        \item Center of second circle: $C_2 = (0, 10)$ (since tangent to $x$-axis at origin with radius 10)
        \item Line through intersection points: call it line $L$
        \item Line through centers: call it line $M$ (radical axis perpendicular)
    \end{itemize}
    \item \textbf{Intuitive Approach}: Use the fundamental property that the line through intersection points of two circles is perpendicular to the line joining their centers (radical axis theorem)
    \item \textbf{Cross-Domain Potential}: 
    \begin{itemize}
        \item Pure synthetic geometry using perpendicularity property
        \item Coordinate geometry by finding circle equations and solving system
        \item Vector approach using perpendicular vectors
    \end{itemize}
\end{itemize}
\end{problemconfigbox}

\begin{strategicbox}
\subsection*{A. Psychological Strategies}
\begin{itemize}[leftmargin=*]
    \item \textbf{Mental Toughness}: This is a standard mid-band problem. If the first approach gets algebraically messy, remember there's likely a cleaner geometric insight. Stay calm and consider perpendicularity properties.
    \item \textbf{Creativity Cultivation}: Consider both synthetic (perpendicular bisector) and analytic (coordinate geometry) approaches. The perpendicular bisector approach is more elegant.
    \item \textbf{Iterative Refinement}: If solving the circle equations directly, verify the algebraic manipulation carefully. If using perpendicularity, verify the slope calculation.
\end{itemize}

\subsection*{B. Investigation Strategies}
\begin{itemize}[leftmargin=*]
    \item \textbf{Startup Orientation}: First, identify the centers of both circles using the tangency conditions. This takes 15-20 seconds and provides the foundation for both solution methods.
    \item \textbf{Get Your Hands Dirty}: Draw a quick sketch showing:
    \begin{itemize}
        \item Circle 1 centered at $(4,0)$ touching $y$-axis at origin
        \item Circle 2 centered at $(0,10)$ touching $x$-axis at origin  
        \item The line through centers
        \item The chord through intersection points (perpendicular to line of centers)
    \end{itemize}
    \item \textbf{Penultimate Step}: The final step is computing $\frac{-1}{m}$ where $m$ is the slope of the line through centers, or simplifying the line equation from the coordinate approach.
    \item \textbf{Make It Easier}: Instead of finding actual intersection points (messy algebra), use the perpendicularity property to avoid solving quadratic systems.
    \item \textbf{Conjecture via Small Cases}: Not directly applicable here, but verify the perpendicular bisector property with simple mental cases.
    \item \textbf{Visualization}: A sketch immediately reveals:
    \begin{itemize}
        \item Centers at $(4,0)$ and $(0,10)$
        \item Line through centers has negative slope
        \item Line through intersections has positive slope (perpendicular)
        \item This eliminates any answer choices with negative slopes or extremely steep slopes
    \end{itemize}
\end{itemize}

\subsection*{C. Argument Construction}
\begin{itemize}[leftmargin=*]
    \item \textbf{Proof Method}: Direct computation using perpendicularity property or algebraic solution
    \item \textbf{Key Lemmas}: 
    \begin{itemize}
        \item The line through intersection points of two circles is perpendicular to the line joining their centers
        \item For perpendicular lines with slopes $m_1$ and $m_2$: $m_1 \cdot m_2 = -1$
    \end{itemize}
    \item \textbf{Logical Structure}: 
    \begin{enumerate}
        \item Determine centers from tangency conditions
        \item Compute slope of line through centers
        \item Apply perpendicularity to find desired slope
    \end{enumerate}
\end{itemize}
\end{strategicbox}

\begin{tacticalbox}
\subsection*{A. Core Mathematical Tactics}
\begin{itemize}[leftmargin=*]
    \item \textbf{Primary Tactics}: 
    \begin{itemize}
        \item \textbf{Perpendicularity Exploitation}: Use the fundamental property that the radical axis (line through intersection points) is perpendicular to the line of centers
        \item \textbf{Symmetry Observation}: The geometric configuration has reflection symmetry considerations that can guide intuition
    \end{itemize}
    \item \textbf{Domain-Specific Tactics}: 
    \begin{itemize}
        \item \textbf{Circle-Line Relationships}: Tangency condition determines center location precisely
        \item \textbf{Radical Axis Theory}: The chord of intersection is the radical axis
        \item \textbf{Slope Relationships}: Perpendicular lines have negative reciprocal slopes
    \end{itemize}
    \item \textbf{Crossover Tactics}: 
    \begin{itemize}
        \item \textbf{Geometric-Algebraic}: Can solve purely geometrically (perpendicularity) or algebraically (solving circle equations)
        \item \textbf{Coordinate Geometry}: Translate geometric tangency into algebraic center coordinates
    \end{itemize}
\end{itemize}
\end{tacticalbox}

\begin{techniquesbox}
\subsection*{A. High-Probability Techniques}
\begin{itemize}[leftmargin=*]
    \item \textbf{Primary Technique}: \textbf{Perpendicular Lines Property} (Radical Axis)
    \item \textbf{Recognition Cues}: 
    \begin{itemize}
        \item Two circles intersecting at two points
        \item Asked for slope of line through intersection points
        \item Centers are easily determined
        \item Classic setup for perpendicular bisector/radical axis
    \end{itemize}
    \item \textbf{Application}: 
    \begin{enumerate}
        \item Find centers: $C_1 = (4, 0)$ and $C_2 = (0, 10)$
        \item Compute slope of line through centers: $m_{centers} = \frac{10-0}{0-4} = -\frac{5}{2}$
        \item Apply perpendicularity: $m_{chord} = -\frac{1}{m_{centers}} = -\frac{1}{-5/2} = \frac{2}{5}$
    \end{enumerate}
\end{itemize}

\subsection*{B. Alternative Techniques}
\begin{itemize}[leftmargin=*]
    \item \textbf{Backup Technique 1}: \textbf{Coordinate Geometry - System of Equations}
    \begin{itemize}
        \item Write circle equations: $(x-4)^2 + y^2 = 16$ and $x^2 + (y-10)^2 = 100$
        \item Expand both equations
        \item Subtract one from the other to eliminate quadratic terms
        \item Result is linear equation representing line through intersections
        \item Extract slope from linear equation
    \end{itemize}
    \item \textbf{Backup Technique 2}: \textbf{Power of a Point / Radical Axis Formula}
    \begin{itemize}
        \item Use radical axis formula directly
        \item For circles with equations $f_1(x,y) = 0$ and $f_2(x,y) = 0$
        \item Radical axis is $f_1(x,y) - f_2(x,y) = 0$
        \item Simplify to get line equation and slope
    \end{itemize}
    \item \textbf{Technique Integration}: 
    \begin{itemize}
        \item Primary method (perpendicularity) is fastest and cleanest
        \item Coordinate method provides verification and alternate path if perpendicularity is forgotten
        \item Both methods should yield $m = \frac{2}{5}$
    \end{itemize}
\end{itemize}
\end{techniquesbox}

\begin{verificationbox}
\subsection*{A. Verification Level Selection}
\begin{itemize}[leftmargin=*]
    \item \textbf{Chosen Level}: Level 3-4 (Targeted Spot Check to Standard Verification)
    \item \textbf{Justification}: 
    \begin{itemize}
        \item Mid-band problem (P10) with standard technique
        \item Straightforward calculation once method is identified
        \item High confidence expected ($\sim$85-90\%)
        \item Need to verify slope calculation and perpendicularity application
    \end{itemize}
    \item \textbf{Time Budget}: 30-45 seconds for verification
\end{itemize}

\subsection*{B. Verification Checklist}
\begin{itemize}[leftmargin=*]
    \item \textbf{Domain Restrictions}: 
    \begin{itemize}
        \item[$\checkmark$] Verify centers are correct: $(4,0)$ and $(0,10)$ from tangency conditions
        \item[$\checkmark$] Confirm circles actually intersect (not disjoint or one inside other)
        \item[$\checkmark$] Distance between centers: $\sqrt{16+100} = \sqrt{116} \approx 10.77$
        \item[$\checkmark$] Sum of radii: $4 + 10 = 14 > 10.77$ (can intersect)
        \item[$\checkmark$] Difference of radii: $|10-4| = 6 < 10.77$ (not one inside other)
    \end{itemize}
    \item \textbf{Computational Accuracy}: 
    \begin{itemize}
        \item[$\checkmark$] Slope of line through centers: $\frac{10-0}{0-4} = \frac{10}{-4} = -\frac{5}{2}$ $\checkmark$
        \item[$\checkmark$] Negative reciprocal: $-\frac{1}{-5/2} = \frac{1 \cdot 2}{5} = \frac{2}{5}$ $\checkmark$
    \end{itemize}
    \item \textbf{Alternative Method}: Expand circle equations and subtract:
    \begin{itemize}
        \item Circle 1: $x^2 - 8x + 16 + y^2 = 16 \Rightarrow x^2 - 8x + y^2 = 0$
        \item Circle 2: $x^2 + y^2 - 20y + 100 = 100 \Rightarrow x^2 + y^2 - 20y = 0$
        \item Subtract: $(x^2 - 8x + y^2) - (x^2 + y^2 - 20y) = 0$
        \item Simplify: $-8x + 20y = 0 \Rightarrow 20y = 8x \Rightarrow y = \frac{8}{20}x = \frac{2}{5}x$
        \item Slope: $m = \frac{2}{5}$ $\checkmark$ (confirms answer)
    \end{itemize}
    \item \textbf{Edge Cases}: 
    \begin{itemize}
        \item[$\checkmark$] Sign check: slope should be positive (perpendicular to negative slope)
        \item[$\checkmark$] Answer format: $\frac{2}{5}$ is already in simplest form
        \item[$\checkmark$] Matches answer choice (E) $\checkmark$
    \end{itemize}
\end{itemize}
\end{verificationbox}

\begin{solutionbox}
\subsection*{A. Step-by-Step Solution}

\subsubsection*{Method 1: Perpendicular Bisector (Recommended - 2 minutes)}

\begin{enumerate}[leftmargin=*]
    \item \textbf{Initial Setup - Determine Centers (30 seconds)}
    
    For the first circle with radius 4, tangent to the $y$-axis at the origin:
    \begin{itemize}
        \item The center must be on the positive $x$-axis
        \item Distance from center to $y$-axis = radius = 4
        \item Therefore, center $C_1 = (4, 0)$
    \end{itemize}
    
    For the second circle with radius 10, tangent to the $x$-axis at the origin:
    \begin{itemize}
        \item The center must be on the positive $y$-axis  
        \item Distance from center to $x$-axis = radius = 10
        \item Therefore, center $C_2 = (0, 10)$
    \end{itemize}
    
    \item \textbf{Strategic Approach - Apply Perpendicularity (30 seconds)}
    
    Key geometric principle: The line through the intersection points of two circles is perpendicular to the line joining their centers.
    
    \item \textbf{Tactical Implementation (45 seconds)}
    
    Compute the slope of the line through the centers $C_1$ and $C_2$:
    \begin{align*}
        m_{centers} &= \frac{y_2 - y_1}{x_2 - x_1} \\
        &= \frac{10 - 0}{0 - 4} \\
        &= \frac{10}{-4} \\
        &= -\frac{5}{2}
    \end{align*}
    
    \item \textbf{Computational Steps (30 seconds)}
    
    Since the line through the intersection points is perpendicular to the line through the centers, we have:
    \begin{align*}
        m_{chord} \cdot m_{centers} &= -1 \\
        m_{chord} \cdot \left(-\frac{5}{2}\right) &= -1 \\
        m_{chord} &= \frac{-1}{-5/2} \\
        m_{chord} &= \frac{2}{5}
    \end{align*}
    
    \item \textbf{Verification (30 seconds)}
    
    Check reasonableness:
    \begin{itemize}
        \item Line through centers has negative slope $\checkmark$
        \item Perpendicular line should have positive slope $\checkmark$
        \item Magnitude check: $\frac{5}{2} \times \frac{2}{5} = 1$ $\checkmark$
    \end{itemize}
    
    \item \textbf{Final Answer}
    
    The slope is $\frac{2}{5}$, which is \textbf{answer choice (E)}.
\end{enumerate}

\subsubsection*{Method 2: Coordinate Geometry (Alternative - 3-4 minutes)}

\begin{enumerate}[leftmargin=*]
    \item \textbf{Write Circle Equations}
    
    Circle 1 centered at $(4, 0)$ with radius 4:
    \begin{equation}
        (x-4)^2 + y^2 = 16
    \end{equation}
    
    Circle 2 centered at $(0, 10)$ with radius 10:
    \begin{equation}
        x^2 + (y-10)^2 = 100
    \end{equation}
    
    \item \textbf{Expand Both Equations}
    
    Expanding equation (1):
    \begin{align*}
        x^2 - 8x + 16 + y^2 &= 16 \\
        x^2 - 8x + y^2 &= 0
    \end{align*}
    
    Expanding equation (2):
    \begin{align*}
        x^2 + y^2 - 20y + 100 &= 100 \\
        x^2 + y^2 - 20y &= 0
    \end{align*}
    
    \item \textbf{Find Radical Axis (Line Through Intersections)}
    
    Subtract the first expanded equation from the second:
    \begin{align*}
        (x^2 + y^2 - 20y) - (x^2 - 8x + y^2) &= 0 \\
        x^2 + y^2 - 20y - x^2 + 8x - y^2 &= 0 \\
        8x - 20y &= 0 \\
        8x &= 20y \\
        y &= \frac{8}{20}x \\
        y &= \frac{2}{5}x
    \end{align*}
    
    \item \textbf{Extract Slope}
    
    The equation $y = \frac{2}{5}x$ is in slope-intercept form, so the slope is $\frac{2}{5}$.
\end{enumerate}

\subsection*{B. Solution Comparison}

\begin{itemize}[leftmargin=*]
    \item \textbf{Method 1 Advantages}: 
    \begin{itemize}
        \item Faster (2 minutes vs 3-4 minutes)
        \item Less algebraic manipulation
        \item Uses elegant geometric insight
        \item Fewer opportunities for calculation error
    \end{itemize}
    \item \textbf{Method 2 Advantages}:
    \begin{itemize}
        \item More mechanical/algorithmic (less insight required)
        \item Provides the actual line equation, not just slope
        \item Can be extended to find exact intersection points if needed
        \item Good fallback if perpendicularity property is forgotten
    \end{itemize}
    \item \textbf{Both methods yield}: $m = \frac{2}{5}$, confirming the answer is \textbf{(E)}.
\end{itemize}
\end{solutionbox}

\begin{competitionbox}
\subsection*{A. Time Management}
\begin{itemize}[leftmargin=*]
    \item \textbf{Estimated Time}: 3-4 minutes total (including verification)
    \begin{itemize}
        \item Setup and center identification: 30-45 seconds
        \item Method 1 (perpendicularity): 1.5-2 minutes
        \item OR Method 2 (coordinate): 2.5-3 minutes
        \item Verification: 30-45 seconds
    \end{itemize}
    \item \textbf{Time Checkpoints}: 
    \begin{itemize}
        \item At 1 minute: Should have centers identified and method chosen
        \item At 2 minutes: Should be computing final slope
        \item At 3 minutes: Should be verifying answer
        \item If stuck after 4 minutes: Flag for review and move on
    \end{itemize}
    \item \textbf{Priority Level}: \textbf{High} - This is an accessible P10 problem with standard techniques. Should be completed before attempting harder problems.
\end{itemize}

\subsection*{B. Risk Assessment}
\begin{itemize}[leftmargin=*]
    \item \textbf{Confidence Level}: 90-95\% (standard problem with well-known techniques)
    \item \textbf{Abandonment Criteria}: 
    \begin{itemize}
        \item If neither perpendicularity nor coordinate methods make progress after 4 minutes
        \item If algebra becomes unmanageably complex (indicates wrong approach)
        \item Should not happen for this problem with proper technique
    \end{itemize}
    \item \textbf{Flag for Review}: No (unless computational error is suspected)
\end{itemize}

\subsection*{C. Answer Choice Strategy}
\begin{itemize}[leftmargin=*]
    \item \textbf{Elimination Potential}:
    \begin{itemize}
        \item Slope should be positive (perpendicular to negative slope) - eliminates any negative options (none present)
        \item Slope should be relatively small (less than 1) since line through centers has slope $-\frac{5}{2}$ (magnitude $> 2$)
        \item This reasoning supports options (A), (B), or (E) over (C) or (D) which involve $\sqrt{29} \approx 5.4$
    \end{itemize}
    \item \textbf{Verification Against Choices}:
    \begin{itemize}
        \item (A) $\frac{2}{7} \approx 0.286$ 
        \item (B) $\frac{3}{7} \approx 0.429$
        \item (C) $\frac{2}{\sqrt{29}} \approx 0.371$
        \item (D) $\frac{1}{\sqrt{29}} \approx 0.186$
        \item (E) $\frac{2}{5} = 0.400$ $\checkmark$
    \end{itemize}
    \item Our answer $\frac{2}{5}$ matches choice (E) exactly.
\end{itemize}
\end{competitionbox}

\begin{keyinsightbox}
\textbf{Key Insight}: The line through the intersection points of two circles is always perpendicular to the line joining their centers. This is a fundamental property of circles (radical axis theorem) that converts a potentially messy coordinate geometry problem into a simple slope calculation.

Once you identify the centers $(4,0)$ and $(0,10)$, you can immediately find the answer using perpendicular slopes without solving any systems of equations or finding actual intersection points.

Formula: If line through centers has slope $m$, then line through intersections has slope $-\frac{1}{m}$.
\end{keyinsightbox}

\begin{warningbox}
\textbf{Warning: Common Pitfalls}

\begin{enumerate}
    \item \textbf{Misidentifying Centers}: Make sure to use the tangency condition correctly. A circle of radius $r$ tangent to an axis at the origin has its center at distance $r$ from that axis.
    
    \item \textbf{Sign Error in Slope}: When computing $\frac{10-0}{0-4}$, don't forget the negative sign. The slope is $-\frac{5}{2}$, not $+\frac{5}{2}$.
    
    \item \textbf{Reciprocal vs Negative Reciprocal}: Perpendicular slopes are \emph{negative reciprocals}, not just reciprocals. The product should be $-1$, not $+1$.
    
    \item \textbf{Unnecessary Complexity}: Don't try to find the actual intersection points unless required. The problem only asks for the slope, which can be found more efficiently using perpendicularity.
    
    \item \textbf{Algebraic Errors in Method 2}: When expanding and subtracting circle equations, track signs carefully. It's easy to drop a negative sign when combining terms.
\end{enumerate}
\end{warningbox}

\begin{tipbox}
\textbf{Pro Tips for Competition Success}

\begin{enumerate}
    \item \textbf{Pattern Recognition}: Whenever you see "two circles intersecting" and need to find something about the line through intersections, think \emph{perpendicular to line of centers} immediately. This is a P9-P12 level standard technique.
    
    \item \textbf{Sketch First}: A quick 15-second sketch reveals the geometric relationship and often suggests the most elegant solution path.
    
    \item \textbf{Time Savings}: Method 1 (perpendicularity) saves 1-2 minutes compared to Method 2 (coordinate algebra). In competition, always look for geometric shortcuts before diving into algebra.
    
    \item \textbf{Verification Efficiency}: Use the alternative method (coordinate) as a quick verification if time permits. Both methods taking 30 seconds each for verification is faster than re-solving completely.
    
    \item \textbf{Answer Choice Sanity Check}: Before detailed calculation, estimate: center line has slope about $-2.5$, so perpendicular should be about $+0.4$. This points to answer (E) and gives confidence in the result.
\end{enumerate}
\end{tipbox}

\section*{Final Answer}

Using the perpendicularity property of circles, we find that the slope of the line through the intersection points is:

\begin{equation*}
\boxed{\textbf{(E) } \frac{2}{5}}
\end{equation*}

\section*{Confidence Assessment}
\begin{itemize}[leftmargin=*]
    \item \textbf{Confidence Level}: 95\% - Standard mid-band problem with well-known geometric property. Verified with two independent methods.
    \item \textbf{Verification Applied}: 
    \begin{itemize}
        \item Level 4 (Standard Verification) 
        \item Primary method: Perpendicular slopes calculation
        \item Cross-check: Coordinate geometry approach yielding same result
        \item Reasonableness: Sign and magnitude checks confirm answer
    \end{itemize}
    \item \textbf{Flag for Review}: No - High confidence with dual-method verification
    \item \textbf{Time Spent}: Approximately 3 minutes for solution + 1 minute for verification = 4 minutes total (within target range for P10)
\end{itemize}

\section*{Pattern Library Entry}

\textbf{Problem Signature}: Circle intersection - find slope of chord

\textbf{Recognition Cues}: 
\begin{itemize}
    \item Two circles with known centers
    \item Asked about line through intersection points
    \item Tangency conditions that determine centers
\end{itemize}

\textbf{Primary Technique}: Perpendicular bisector property (radical axis)

\textbf{Key Formula}: If centers form line with slope $m$, intersection chord has slope $-\frac{1}{m}$

\textbf{Alternative Approaches}: Coordinate geometry (expand circles, subtract equations)

\textbf{Common Traps}: Sign errors in slope calculation, forgetting negative reciprocal

\textbf{Time Estimate}: 3-4 minutes

\textbf{Similar Problems}: AMC12 problems involving radical axis, perpendicular bisectors of chords, circle geometry with coordinate systems

\end{document}

